\documentclass[11pt,a4paper]{article}

% --- Packages ---
\usepackage[margin=.75in]{geometry}
\usepackage[utf8]{inputenc}
\usepackage[T1]{fontenc}
\usepackage{lmodern} % Fix for font not found error
\usepackage{titlesec}
\usepackage{enumitem}
\usepackage{hyperref}
\usepackage{xcolor}
\usepackage{fancyhdr}
\usepackage{comment}
\usepackage{cleveref}
\usepackage[bottom]{footmisc}
\usepackage{graphicx}
\usepackage{longtable}
\usepackage{booktabs}
\usepackage{array}
\usepackage{calc}
\usepackage{etoolbox}

% Standard Pandoc setup
\providecommand{\tightlist}{%
  \setlength{\itemsep}{0pt}\setlength{\parskip}{0pt}}

\crefformat{section}{\S#2#1#3}
\crefformat{subsection}{\S#2#1#3}
\crefformat{subsubsection}{\S#2#1#3}
\crefformat{figure}{#2Figure~#1#3}
\crefformat{footnote}{#2\footnotemark[#1]#3}

\linespread{0.9} % Riduce lo spazio tra le linee

% --- Colors ---
\definecolor{primary}{RGB}{0, 0, 0}
\definecolor{links}{HTML}{0000EE}
\definecolor{blue}{HTML}{26428B}

\hypersetup{
    colorlinks=true,
    linkcolor=links,
    urlcolor=links,
    citecolor=links,
}

% --- Section Formatting ---
\titleformat{\section}
  {\normalfont\Large\bfseries\color{blue}}{\thesection}{1em}{}
\titleformat{\subsection}
{\normalfont\large\bfseries\color{blue}}{\thesubsection}{1em}{}
\titleformat{\subsubsection}
{\normalfont\bfseries\color{blue}}{\thesubsubsection}{1em}{}

% --- Header and Footer ---
\pagestyle{fancy}
\fancyhead[L]{\small Giuseppe Capasso}
\fancyhead[R]{\small intro-to-linux}
\fancyfoot[C]{\thepage}
\renewcommand{\headrulewidth}{0.4pt}

% --- Custom Title Command ---
\newcommand{\proposaltitle}[1]{
    \begin{center}
        \vspace*{0.1cm}
        {\LARGE \bfseries \color{blue} #1} \\[1.5ex]
        {\large \textbf{Giuseppe Capasso}} \\[1ex]
                \small{
            \vspace{0.1cm}
            \textbf{Materia:} intro-to-linux\\
        }
            \end{center}
    \vspace{0.3cm}
    \hrule height 0.5pt
    \vspace{0.5cm}
}

\begin{document}

\proposaltitle{Dispensa ITS: Storage Management & Backup}

Dispensa ITS: Storage Management \& Backup Strategies Modulo: Linux
System Administration - Lezione 3 Obiettivo: Gestire dischi, partizioni,
file system e implementare strategie di backup sicure.

\begin{enumerate}
\def\labelenumi{\arabic{enumi}.}
\tightlist
\item
  Architettura dello Storage in Linux A differenza di Windows, che
  assegna lettere di unità (C:, D:), Linux gestisce i dispositivi di
  archiviazione come file speciali situati nella directory /dev. +1
\end{enumerate}

1.1 Naming Convention (Nomenclatura) Il Kernel assegna i nomi in base al
tipo di interfaccia: Nome Device /dev/sdX /dev/sdXN

Descrizione SATA/SCSI/USB Device. La `X' è una lettera progressiva (a,
b, c\ldots). Partizione. La `N' è un numero.

/dev/ NVMe SSD. Dischi moderni ad alta velocità. nvme0n1 /dev/sr0 CD/DVD
ROM. Esporta in Fogli

Esempio /dev/sda (1° disco), /dev/sdb (2° disco) /dev/sda1 (1ª
partizione del 1° disco) /dev/nvme0n1p1 (Partizione 1) -

\begin{enumerate}
\def\labelenumi{\arabic{enumi}.}
\setcounter{enumi}{1}
\tightlist
\item
  Guida Pratica: Aggiungere un Disco in VirtualBox Prima di poter
  partizionare un disco, dobbiamo collegarlo fisicamente (o
  virtualmente) alla macchina. Ecco la procedura passo-passo per
  simulare l'inserimento di un secondo Hard Disk (HDD) nella vostra VM
  Ubuntu. Procedura (A VM Spenta):
\item
  Apri VirtualBox e seleziona la tua VM Ubuntu dalla lista a sinistra.
\item
  Clicca su Impostazioni (Settings) -\textgreater{} Archiviazione
  (Storage).
\item
  Nell'albero ``Dispositivi di archiviazione'', seleziona il Controller:
  SATA.
\item
  Clicca sull'icona ``Aggiunge disco rigido'' (quella con il + verde
  sopra l'icona del disco, non quella del CD).
\end{enumerate}

5. Nella finestra che si apre, clicca su Crea (Create). 6. Tipo di file:
Seleziona VDI (VirtualBox Disk Image) -\textgreater{} Avanti. 7.
Archiviazione: Seleziona Allocato dinamicamente (Occupa spazio solo se
lo riempi) -\textgreater{} Avanti. 8. Nome e Dimensione: • Nome:
Disco\_Dati (o simile). • Dimensione: 5 GB (sufficienti per
l'esercizio). 9. Clicca su Fine. 10.Ora vedrai il nuovo disco nella
lista ``Not Attached''. Selezionalo e clicca su Scegli (Choose).
11.Clicca OK per chiudere le impostazioni e avvia la VM. Verifica in
Linux: Una volta avviata la VM, apri il terminale e digita: Bash lsblk

Dovresti vedere un nuovo dispositivo (probabilmente sdb) di circa 5GB
senza partizioni sotto di esso.

\begin{enumerate}
\def\labelenumi{\arabic{enumi}.}
\setcounter{enumi}{2}
\tightlist
\item
  Partizionamento e File System Un disco grezzo non è utilizzabile. Deve
  essere partizionato e formattato.
\end{enumerate}

3.1 MBR vs GPT Esistono due standard per la tabella delle partizioni:
Caratteristica MBR (Master Boot Record) GPT (GUID Partition Table)
Limite Dimensione Max 2 TB Fino a 9.4 Zettabyte (praticamente
illimitato) Partizioni Max 4 Primarie Fino a 128 partizioni Robustezza
Singola copia all'inizio del disco Copia di backup alla fine del disco
Utilizzo Legacy / BIOS vecchi Standard moderno / UEFI Esporta in Fogli

3.2 Strumenti di Partizionamento Utilizzeremo fdisk (più comune per
script/MBR) o cfdisk (interfaccia grafica testuale, più semplice).
Esempio con fdisk (su /dev/sdb): Bash

sudo fdisk /dev/sdb

• m: Mostra il menu di aiuto. • n: Nuova partizione. • p: Primaria. • w:
Scrivi le modifiche ed esci.

3.3 Formattazione (Creazione File System) Dopo aver creato la partizione
(es. /dev/sdb1), dobbiamo decidere ``come'' scriverci i dati. Linux usa
principalmente ext4 (Fourth Extended Filesystem). Comando per
formattare: Bash sudo mkfs.ext4 /dev/sdb1

Attenzione: mkfs cancella irreversibilmente tutti i dati nella
partizione specificata!

\begin{enumerate}
\def\labelenumi{\arabic{enumi}.}
\setcounter{enumi}{3}
\tightlist
\item
  Mounting: Accedere ai Dati In Linux, per accedere a un disco
  formattato, devi ``montarlo'' su una cartella esistente.
\end{enumerate}

4.1 Mount Temporaneo 1. Crea il punto di mount (una cartella vuota):
Bash sudo mkdir /mnt/dati

\begin{enumerate}
\def\labelenumi{\arabic{enumi}.}
\setcounter{enumi}{1}
\item
  Collega la partizione alla cartella: Bash sudo mount /dev/sdb1
  /mnt/dati
\item
  Verifica: Bash df -h
\end{enumerate}

Nota: Al riavvio, questo collegamento andrà perso.

4.2 Mount Permanente (/etc/fstab) Per rendere il mount automatico
all'avvio, dobbiamo modificare il file /etc/fstab.

1. Ottieni l'UUID (codice univoco) del disco (è più sicuro usare l'UUID
che il nome /dev/sdb1, perché i nomi possono cambiare se sposti i cavi):
Bash sudo blkid

Copia la stringa UUID relativa a /dev/sdb1. 2. Modifica il file fstab:
Bash sudo nano /etc/fstab

\begin{enumerate}
\def\labelenumi{\arabic{enumi}.}
\setcounter{enumi}{2}
\tightlist
\item
  Aggiungi una riga alla fine con questa sintassi: {[}Device{]} {[}Mount
  Point{]} {[}Filesystem{]} {[}Options{]} {[}Dump{]} {[}Pass{]} Esempio:
  Plaintext UUID=1234-5678-abcd
\end{enumerate}

/mnt/dati

ext4

defaults

0

2

ATTENZIONE: Un errore di sintassi in /etc/fstab può impedire l'avvio del
sistema. Verificare sempre con sudo mount -a prima di riavviare.

\begin{enumerate}
\def\labelenumi{\arabic{enumi}.}
\setcounter{enumi}{4}
\tightlist
\item
  Strategie di Backup Il backup non è solo ``copiare file'', ma
  pianificare il recupero (Restore). 5.1 Tipologie di Backup Tipo
\end{enumerate}

Descrizione Copia completa di tutti i Full dati selezionati. Copia solo
i dati modificati dall'ultimo Incrementale backup (di qualsiasi tipo).
Copia i dati modificati Differenziale dall'ultimo Full Backup. Esporta
in Fogli

Pro (Vantaggi) Contro (Svantaggi) Restore velocissimo (basta Lento da
eseguire, occupa l'ultimo file). tantissimo spazio. Restore lento e
complesso (serve Backup rapidissimo, poco il Full + tutti gli
incrementali spazio occupato. successivi). Via di mezzo. Restore più
Occupa progressivamente più veloce dell'incrementale. spazio
dell'incrementale.

5.2 La Regola del 3-2-1 La strategia standard industriale per la
sicurezza dei dati: 1. Mantenere 3 copie dei dati (1 originale + 2
backup).

2. Salvare su 2 tipi di supporto diversi (es. Disco locale + NAS, oppure
Disco + Cloud). 3. Conservare 1 copia Off-site (fuori sede, per
proteggersi da incendi/furti fisici).

\begin{enumerate}
\def\labelenumi{\arabic{enumi}.}
\setcounter{enumi}{5}
\tightlist
\item
  Strumenti di Backup Linux 6.1 Tar (Tape Archive) Storico strumento per
  impacchettare file. Utile per backup completi (Full) o archiviazione
  ``a freddo''. • Creare un archivio compresso (gzip): Bash tar -czvf
  backup\_home.tar.gz /home/studente
\end{enumerate}

• -c: Create • -z: Gzip compression • -v: Verbose (mostra i file) • -f:
File (nome del file di output) • Estrarre (Restore): Bash tar -xzvf
backup\_home.tar.gz -C /tmp/restore

6.2 Rsync (Remote Sync) Lo strumento più potente per backup incrementali
e sincronizzazione. Copia solo i ``delta'' (le differenze), risparmiando
banda e tempo. • Sintassi base: Bash rsync -av /sorgente /destinazione

• -a (Archive): Mantiene permessi, date, link simbolici, utenti
(ricorsivo). • -v (Verbose): Mostra dettagli. • Backup speculare
(Mirroring): Bash rsync -av --delete /home/studente/progetti
/mnt/backup/

• --delete: Attenzione! Cancella nella destinazione i file che non
esistono più nella sorgente. Crea una copia esatta 1:1.



\end{document}
